\phantomsection
\chapter*{Hope}
\addcontentsline{toc}{section}{Hope --- Khalqiyah Ahmad Siddiq}
\penulis{Khalqiyah Ahmad Siddiq}

\awalcerpen{S}{aat dunia berada di tepi kehancuran}, sebuah kata sering diteriakkan oleh mereka yang telah kehilangan asa: harapan. Akan tetapi, bagi sebagian orang lain, kata tersebut hanyalah sekadar omong kosong belaka. Apabila harapan itu memang sungguh ada, lalu di manakah ia berada ketika peradaban ini rata dengan tanah?

Sudah tujuh puluh enam tahun berlalu sejak bencana besar merenggut wajah bumi yang dahulu hijau. Dunia tidak pernah lagi sama seperti dulu. Mayoritas umat manusia tewas dan tersapu oleh keganasan alam, sementara segelintir yang bertahan memilih merangkak dan meringkuk di bungker-bungker bawah tanah untuk mempertahankan hidup. Banyak yang menduga bahwa malapetaka ini disebabkan oleh ledakan senjata pemusnah massal. Tetapi, dugaan tersebut sangatlah keliru. Kiamat yang berlangsung perlahan ini muncul akibat pemanasan global yang sangat ekstrem, memicu lahirnya badai debu abadi yang menutupi atmosfer tanpa henti.

Di bawah naungan tanah yang pengap, bertahan hidup merupakan perjuangan yang penuh siksaan. Persediaan air bersih telah habis menguap, ditelan panas bumi yang tak tertahankan. Agar napas tetap berhembus, manusia terpaksa mendaur ulang cairan dari keringat dan urin mereka sendiri. Setiap tetes yang melewati kerongkongan terasa begitu menyiksa, menjadi penanda betapa ironisnya kehidupan saat ini.

Di koridor Bungker Sektor B0910 yang remang-remang, suara batuk berat memecah kesunyian yang mencekam.

``Milen$\dots$ uhuk! uhuk! Milen$\dots$''

Panggilan itu terdengar serak dan rapuh dari balik bilik yang berdebu.

``Ada apa, Bu?'' Milen memanggil dengan cepat.

Seorang perempuan tergeletak lemah di atas pembaringan yang sudah usang. Wajahnya tampak pucat sekali, namun ia masih berupaya memperlihatkan ketegaran walau tubuhnya tengah dilahap penyakit yang sangat berat.

``Datanglah sebentar, Nak,'' bisiknya lirih.

Milen melangkah mendekat dengan dada yang dicekam kecemasan yang luar biasa. ``Ada apa, Bu?''

Sambil menahan getaran dan deru batuk yang tak henti menyiksa dadanya, ibu itu menatap mata anaknya dengan penuh kelembutan. ``Ibu sudah tak sanggup lagi, Milen$\dots$ Kamu harus mandiri dan tegar menghadapi apa pun di dunia ini, sebab harapan akan selalu menyertaimu. Tolong jaga adikmu dengan baik, didiklah dia seperti cara Ibu mendi$\dots$''

Kalimat itu terhenti. Napas yang tersisa mengembus perlahan, meninggalkan jemari yang terkulai lemas di atas pangkuan.

``Ibu? Ibu?'' Tangis Milen meledak memecah keheningan ruangan. Ia mengguncang perlahan tubuh yang semakin dingin itu dengan dada yang sesak oleh duka. ``Ibu tidak apa-apa, kan? Ibu$\dots$ jawab, Bu! Ibu!''

Sudah sepuluh tahun berlalu sejak malam yang kelam itu. Sekarang aku sudah menginjak usia sembilan belas tahun, sedangkan adikku, Reymi, baru saja genap dua belas tahun. Mau tak mau, kami bergantung pada paman kami—satu-satunya adik kandung almarhumah Ibu. Akan tetapi, Paman sungguh jauh berbeda dari Ibu. Di matanya, alkohol adalah candu, dan kami hanya dianggap hama yang mengganggu pandangan. Tiap kali malam datang dan bau racun pemabuk menguasai kepalanya, siksaan demi siksaan menjadi rutinitas yang harus kami terima.

Puncaknya adalah saat Paman memukuli Reymi tanpa ampun, meninggalkan luka robek yang membekas parah di wajah mungil adikku. Amarah membakar dadaku, namun rasa takut membelenggu jemariku untuk melawan. Malam itu, aku sadar kami tidak bisa bertahan lebih lama lagi. Kami harus meninggalkan sektor ini, memilih jalan kami sendiri agar tetap hidup.

Sekitar pukul satu dini hari, ketika seluruh sektor diliputi kegelapan, aku dan Reymi menyelinap keluar dari kamar dengan napas tertahan. Suara tawa sumbang dan celotehan Paman terdengar dari arah ruang tamu. Dia sedang mabuk berat, tertawa menjijikkan bersama seorang wanita asing yang tidak kami kenal.

Kami sempat berhenti sejenak di balik gelapnya dinding. Tepat saat Paman hendak menarik wanita itu ke dalam pelukannya, nasib buruk menimpa kami. Langkah Reymi tersandung. Tubuh kecilnya jatuh ke lantai, menimbulkan bunyi dentuman pelan yang bergema tajam di tengah keheningan.

``Heeeey! Siapa itu?!'' teriak Paman, suaranya parau dan penuh ancaman.

Rasa panik menyelimuti diriku. Aku segera memeluk tubuh Reymi untuk menolongnya berdiri, namun Paman bergerak jauh lebih cepat. Dengan kejam, ia mencengkeram kerah baju Reymi, menyeret adikku yang menjerit ketakutan ke ruangan sebelah, kemudian mengikatnya tanpa sedikit pun belas kasihan.

``Bocah sialan! Kenapa kamu selalu mengganggu, hah?! Tidak bisakah kamu membiarkanku tenang walau sebentar saja?!'' hardik Paman dengan mata merah membara.

``Paman, lepaskan dia!'' teriakku histeris, berusaha menerobos masuk.

Namun, kegilaan malam itu berakhir menjadi mimpi buruk yang sama sekali tak pernah kubayangkan. Di bawah pengaruh alkohol yang menguasai kesadarannya, gerakan Paman menjadi liar dan tak terkendali. Pukulan serta dorongan kasarnya menghantam tubuh mungil Reymi dengan terlalu keras. Dalam hitungan detik yang menegangkan, napas adikku berhenti. Tubuh kecil itu terkulai dingin di atas lantai, tak bergerak lagi.

Duniaku runtuh seketika. Almarhumah Ibu mempercayakan Reymi kepadaku, dan kini adikku tewas di depan mataku sendiri.

``BAJINGAN! KENAPA KAU MEMBUNUHNYA?!'' jeritku melengking, dipenuhi erangan murka dan dendam yang membakar habis seluruh rasa takutku.

Paman menoleh dengan lambat, matanya menyalang penuh kebencian. ``Bocah kurang ajar! Beraninya kamu menyebutku bajingan?!''

Ia langsung menyambar tongkat \textit{baseball} berbahan besi yang tersandar di sudut ruangan. Tanpa aba-aba, Paman mengayunkannya dengan sangat brutal. Aku mencoba menerjang untuk membalas, namun ayunan keras itu lebih dulu menghantamku.

Brak!

Hantaman keras itu menghantam sisi kepalaku. Pandanganku langsung memutih, dan rasa sakit yang tak terkira menusuk sampai ke dasar tengkorak. Darah yang hangat dan kental segera mengalir membasahi pelipis hingga menutupi sebagian pandanganku. Aku jatuh tersungkur di lantai yang dingin sambil mengerang kesakitan, sementara Paman kembali mengangkat tongkatnya dengan seringai setan.

Karena sadar tak mungkin menang dalam keadaan seperti itu, aku memaksa tubuhku yang gemetar untuk berdiri. Dengan sisa tenaga dan darah yang masih mengucur, aku berbalik lalu berlari menembus gelapnya malam.

Aku berlari tanpa tujuan, meraba dinding-dinding dingin koridor bungker itu, meninggalkan jasad adikku and rumah penuh neraka tersebut. Air mata yang bercampur darah mengalir deras di pipiku. Di tengah tangis dan rasa sesak yang menghancurkan dada, satu tekad membara dalam jiwaku: aku akan bertahan hidup, dan suatu hari nanti, aku akan kembali untuk menuntut balas.

Dengan kepala berlumuran darah dan kesadaran yang semakin menipis, aku terus memaksa kakiku untuk melangkah. Sektor E457—tempatku berada sekarang—terbilang sangat jauh dari Sektor B yang baru saja kutinggalkan. Rasa sakit yang luar biasa menusuk tajam di kepalaku, menyiksaku di setiap langkah. Akhirnya, tubuhku sampai pada batas kemampuannya. Sebelum pandanganku benar-benar hilang menjadi kegelapan total dan tubuhku jatuh ke tanah, aku sempat melihat sesosok kakek tua yang berdiri bersandar pada sebuah tongkat.

Kemudian, semuanya berubah menjadi gelap.

Saat mataku terbuka kembali, aku mendapati diriku sedang terbaring di sebuah ruangan yang asing. Dinding-dinding ruangan itu penuh dengan peralatan canggih serta panel teknologi yang belum pernah kulihat sebelumnya.

``Di mana aku berada? Tempat apa ini?'' bisikku pelan, sambil memegangi kepala yang berdenyut kencang, berupaya menyusun kembali ingatan tentang tragedi berdarah yang baru saja berlangsung.

Dari kejauhan, terdengar suara langkah kaki yang semakin mendekat. Sosok kakek yang kulihat sebelum pingsan kini berjalan menghampiriku. Wajahnya terlihat teduh ketika ia berkata, ``Tenanglah, Nak. Di sini aman. Tadi malam aku menemukanmu tergeletak dengan kepala berlumuran darah, jadi aku membawamu ke rumahku.''

Aku menatapnya dengan bingung dan waspada. ``Anda$\dots$ siapa?''

Kakek tua itu tersenyum tipis dengan sikap yang sangat elegan. ``Ah, benar. Aku belum memperkenalkan diri—''

``Kakek, apakah dia sudah siuman?''

Ucapan sang kakek terputus oleh suara nyaring seorang gadis yang tiba-tiba masuk melangkah.

``Sudah, dia baru saja bangun,'' jawab kakek itu dengan lembut.

Gadis itu melangkah mendekati tempat tidurku. Kakek itu lalu meneruskan perkenalannya, ``Gadis yang ada di sampingmu ini adalah Kiren, cucuku. Dan namaku Dokter Heri.''

Jantungku langsung berdetak kencang. Nama itu terdengar familier. Wajahnya pernah muncul di halaman koran lama yang kerap kubaca—dia seorang ilmuwan jenius yang mengabdikan hidupnya untuk mencari cara mengembalikan kondisi bumi seperti semula.

Menyadari perubahan pada raut wajahku, Dokter Heri memperhatikanku dengan saksama. ``Hmm$\dots$ dari ekspresimu, sepertinya kamu mengenalku, ya?''

``E-eh, benar, Dok,'' jawabku gugup sambil menelan ludah. ``Saya pernah melihat foto Dokter di koran.''

Dokter Heri tertawa kecil. ``Hahaha, ternyata saya memang cukup terkenal. Tapi ngomong-ngomong, apa yang terjadi dengan kepalamu? Lukanya sangat parah kemarin.''

Aku terdiam sesaat. Aku tidak ingin orang lain mengetahui tragedi mengerikan di rumah paman, terlebih soal kematian Reymi. Aku memutuskan untuk berbohong.

``Ini$\dots$ hanya akibat perkelahian jalanan, Dok,'' kilahku pelan.

Kiren melipat kedua tangan di depan dada, memandangku dengan tatapan tidak percaya. ``Apa kamu memang punya hobi mencari masalah, ya?''

``I$\dots$ iya, hehe,'' jawabku semakin gelisah, lalu mengalihkan pandangan ke arah lain.

Gadis itu menggelengkan kepala sambil memasang raut wajah keheranan. ``Kenapa, sih, laki-laki itu selalu saja suka membuat masalah?''

``Sudah, sudah, Kiren,'' sela Dokter Heri menenangkan cucunya, kemudian kembali menatapku dengan hangat. ``Oh iya, kita belum berkenalan. Siapa namamu, Nak?''

Aku menghela napas panjang, berusaha menguatkan diri. ``Namaku$\dots$ Milen.''

Beberapa hari pun berlalu. Kondisi fisikku semakin membaik dari hari ke hari. Pada suatu pagi, Dokter Heri mendatangiku dan menyodorkan sebuah tawaran yang sama sekali tidak pernah kubayangkan sebelumnya.

``Apakah kamu mau bergabung dengan organisasiku, Milen?'' tanyanya sambil menatapku lekat.

Aku mengerutkan dahi, lalu tertawa kecil karena kebingungan. ``Hah? Organisasi apa, Dok?''

Dokter Heri menjelaskan bahwa organisasi yang ia pimpin bergerak dalam misi pencarian bahan dan komponen khusus untuk mewujudkan sebuah teknologi besar yang sedang ia kembangkan.

``Kamu akan dikirim ke luar bungker untuk mencari sisa-isa teknologi manusia yang terlupakan dan ditinggalkan di balik tembok sana,'' ujarnya lagi. ``Di setiap misi, kamu akan bertugas bersama dua orang mitra.''

Aku memandangnya dengan ragu. ``Kenapa Dokter memilihku?''

Dokter Heri tersenyum, raut wajahnya menunjukkan rasa kagum. ``Kamu punya ketahanan tubuh yang luar biasa. Ketika kepalamu terluka parah kemarin, kamu masih mampu berjalan sejauh itu. Padahal menurut hasil pemindaian komputer, kondisimu sangat kritis. Selain itu, proses penyembuhan jaringan tubuhmu berlangsung sangat cepat. Kemampuan ini benar-benar cocok untuk tugas di luar bungker yang penuh badai ganas.''

Aku terdiam sesaat, mencerna semua ucapannya. ``Beri aku waktu untuk memikirkannya dulu, Dok.''

Setelah Dokter Heri pergi, aku berjalan mengelilingi area itu. Tempat ini terasa sangat asing bagiku. Di tiap sudut ruangan, terpasang aneka perangkat canggih yang tidak akan pernah ditemukan di bungker pemukiman biasa.

Saat menyusuri koridor yang panjang, mataku menangkap pemandangan menakjubkan dari kejauhan—sebuah taman buatan yang sungguh indah. Di sana, kulihat Kiren duduk sendirian, menikmati suasana tenang di antara pepohonan hijau.

Aku melangkah mendekat, berusaha mencairkan suasana dengan obrolan ringan. ``Kenapa kamu duduk di sini?''

Kiren menoleh perlahan, menyunggingkan senyum tipis. ``Sebab ini adalah tempat favoritku dan Ibu pada masa lalu.''

Rasa penasaran itulah yang mendorong keberanianku untuk bertanya lebih dalam. ``Sekarang ibumu ada di mana?''

Seketika itu juga, wajah Kiren berubah muram. Kepalanya ia tundukkan dengan sangat dalam. ``Dahulu Ibu turut serta bersama Ayah dalam ekspedisi menuju permukaan bumi. Namun di tengah perjalanan, badai pasir raksasa tiba-tiba menerjang dan menewaskan seluruh anggota tim ekspedisi$\dots$ Hanya Kakek yang menjadi satu-satunya orang yang selamat.''

Dada ini terasa berdebar-debar. Perasaan bersalah pun langsung menghantuiku. ``Oh$\dots$ maafkan aku, aku tidak mengetahuinya.''

Kiren kembali mengangkat wajahnya, menyunggingkan senyum tulus walau matanya berkaca-kaca. ``Tidak masalah, lagipula peristiwa itu sudah lama terjadi. Ngomong-ngomong, orang tuamu sekarang di mana? Apa mereka tidak mencarimu?''

Aku terdiam sejenak. Kenangan pahit itu kembali merayap dalam pikiranku. ``Sebenarnya$\dots$ Ibuku telah meninggal sepuluh tahun silam, ketika usiaku baru sembilan tahun, akibat sakit. Dan Ayahku$\dots$ dia telah wafat sebelum aku dilahirkan.''

Kiren menatapku dengan lembut, mengangguk pelan. ``Ternyata$\dots$ nasib kita tak jauh berbeda, ya.''

Setelah cukup lama berbincang dengan Kiren, aku pun kembali ke kamarku. Dalam keheningan ruangan itu, bayang-bayang masa lalu dan bisikan dendam kembali bergaung di kepalaku. Tetapi, perkataan Kiren tentang impiannya menyaksikan bumi yang indah seperti dahulu kala tiba-tiba terlintas. Harapan sederhana itu menyalakan dorongan baru di dadaku—motivasi kuat untuk benar-benar menginjakkan kaki di permukaan atas. Sosoknya yang tegar walau menanggung duka yang sama membuatku merasa tak lagi sendiri di dunia yang rusak ini.

Esok harinya, aku mantap menemui Dokter Heri dan menerima tawaran ekspedisi itu. Di sana, aku diperkenalkan kepada dua anggota tim yang akan menemaniku melangkah ke atas. Yang pertama adalah Horus, pria berbadan kekar yang memancarkan kekuatan fisik luar biasa. Yang kedua adalah Ciberus, sosok pendiam dan dingin dengan tatapan mata yang sungguh mengerikan—jelas menandakan bahwa ia bukan orang sembarangan. Misi ekspedisi kami baru akan dijalankan satu tahun lagi, menunggu badai ekstrem di permukaan mereda.

Karena belum ada tugas resmi, aku memutuskan untuk mencari Kiren guna menghabiskan waktu luang.

``Hei, Kiren! Apakah kamu punya waktu?'' seruku dari kejauhan.

Kiren mengernyit sambil melambaikan tangan. ``Apa? Tidak terdengar!''

Aku berjalan mendekat dan mengulanginya, ``Apakah kamu punya waktu?''

``Memangnya kenapa?'' tanya Kiren sambil memiringkan kepala.

``Aku ingin mengajakmu keluar untuk mencari makanan,'' ajakku.

Kiren terdiam sejenak lalu tersenyum. ``Hmm$\dots$ boleh juga, aku lapar.''

Kami pun menyusuri lorong bungker untuk mencari tempat makan. Sebagian besar kedai terlihat tutup karena krisis pangan dan air bersih yang masih berlangsung. Untunglah, mataku menangkap satu restoran yang masih buka.

``Akhirnya kita menemukan tempat yang buka!'' seruku dengan lega.

Kami masuk, menikmati hidangan yang tersedia, and menghabiskan waktu sambil bertukar cerita.

Lima bulan berlalu begitu cepat. Sekarang, hariku sepenuhnya terisi oleh latihan berat untuk persiapan ekspedisi tujuh bulan ke depan. Sebagian besar waktuku digunakan untuk mengasah kemampuan bertarung tangan kosong.

Ketika aku sedang bermandikan keringat di tengah sesi latihan, Kiren menghampiri sambil melempar sebotol jatah air ke arahku.

``Kamu mau membuat dirimu pingsan sebelum misinya dimulai?'' Pandangannya kosong, kedua tangan disilangkan di depan dada.

Kutangkap botol itu, lalu kuseka keringat di dahi sambil tersenyum tipis. ``Latihan inilah yang akan membuatku tetap bertahan hidup di atas sana.''

Kiren mendengus lirih sambil menatapku dengan tatapan lurus. ``Meski begitu, istirahat bukanlah sebuah pilihan, Milen. Jangan sampai kau jatuh hanya karena lupa mengatur napas.''

Walau nada suaranya terdengar ketus, matanya memperlihatkan kekhawatiran yang sungguh nyata. Di dunia yang kejam ini, kepedulian semacam itu merupakan hal langka yang diam-diam sangat kuhargai.

Sebulan sebelum misi, intensitas latihanku kian bertambah. Malam itu, di koridor lengang dekat area pelatihan, Kiren duduk menantiku. Cahaya yang redup membuat keheningan di antara kami terasa lebih damai daripada biasanya.

``Ekspedisi tinggal sebulan lagi,'' katanya pelan, memecah keheningan sambil menatap lantai.

``Iya. Persiapannya hampir rampung,'' jawabku, lalu duduk di sampingnya dengan napas yang masih belum teratur.

Kiren menoleh, menatapku lekat-lekat tanpa sedikit pun kepalsuan. ``Milen$\dots$ berjanjilah kau akan kembali dengan selamat. Dunia ini sudah merenggut terlalu banyak orang dariku, dan aku$\dots$ tak ingin kehilanganmu juga.''

Melihat tatapannya yang begitu rapuh namun tulus, aku paham apa yang tersimpan di hatinya—perasaan yang sebenarnya juga tumbuh di dadaku. Akan tetapi, karena menyadari betapa berisiko misi yang menanti di depanku, aku memutuskan untuk menahan diri.

``Aku akan kembali, Kiren,'' kataku pelan tetapi tegas sambil menatap kedua matanya. ``Setelah semua kekacauan ini berakhir, kita akan membicarakan masa depan kita. Aku berjanji.''

Kiren terdiam sesaat, kemudian mengangguk pelan dengan senyum tipis yang penuh kelegaan, seakan janji singkat itu sudah cukup memberinya ketenangan.

Sampai tibalah sehari sebelum keberangkatan misi. Aku menjalani rutinitas seperti biasanya, menempa tubuh dan mengasah kemampuan bertarung. Tetapi, keheningan latihan tiba-tiba terpecah. Dari kejauhan, jeritan histeris Kiren menggema, disusul ledakan dahsyat yang mengguncang fondasi fasilitas Dokter Heri.

Jantungku seperti berhenti berdetak. Aku berlari secepat kilat menuju sumber suara.

Pemandangan di hadapanku berubah menjadi neraka. Kobaran api menjulang tinggi, menerangi kegelapan fasilitas, sementara asap hitam mengepul tebal membawa aroma hangus yang mencekam. Di antara remang lidah api, kulihat bayangan dua sosok yang sangat kukenal: Horus dan Ciberus.

Langkahku langsung terhenti oleh kengerian di depan mata. Dokter Heri telah terkapar bersimbah darah, tak bernyawa. Dan sebelum akal sehatku mampu mencerna tragedi itu, kedua pengkhianat tersebut berbalik menyerang Kiren.

``TIDAAAK!'' jeritku, menerjang maju menembus kobaran api dengan sisa tenaga yang kumiliki.

Namun, jarak dan waktu bersekongkol melawanku. Di hadapan mataku sendiri, tebasan serta pukulan mematikan dari mereka merampas tubuh kecil Kiren. Ia jatuh tersungkur ke lantai dingin, sementara Horus dan Ciberus lenyap di balik gumpalan asap tebal.

Aku berlutut di sisinya, mendekap kepala Kiren dengan jari-jari yang bergetar hebat. Darah hangat membasahi tanganku, sama persis seperti malam ketika adikku dibunuh.

Dengan kesadaran yang semakin memudar dan napas yang tersengal, Kiren menatapku. Matanya yang biasanya berkilau kini meredup dengan cepat. ``Cepat$\dots$ pergilah dari sini, Milen$\dots$'' bisiknya sangat pelan, terpotong-potong oleh darah yang membasahi bibirnya. ``M-mereka memburumu$\dots$ Mereka adalah agen dari \textit{Dark Society}$\dots$ Mereka hendak membuat misi kita gagal dan menyabotase semuanya$\dots$ Milen$\dots$ kamu$\dots$ harus tetap hidup$\dots$ Dan kamu$\dots$ adalah harapan umat manusia$\dots$''

Dada Kiren mengembuskan napas terakhirnya. Lengannya terkulai lemas, dan kehangatan tubuhnya perlahan menguap ditelan dinginnya malam.

Duniaku runtuh dalam sekejap. Remuk menjadi serpihan kecil sampai tidak ada yang tersisa.

Seakan-akan dunia ini punya rencana keji yang tak pernah mengizinkanku merasakan kebahagiaan. Pertama-tama, ibuku diambil oleh takdir. Setelah itu, adikku Reymi dibunuh di depan mataku sendiri. Dan sekarang, Kiren—satu-satunya sandaran tempatku menaruh sisa kemanusiaanku—direnggut paksa dari pelukanku. Setiap kali aku berusaha mencintai, setiap kali aku mencoba menyusun kembali benang harapan, dunia akan dengan serakah menelannya bulat-bulat, lalu meludahkan rasa sakit itu tepat ke wajahku.

Di tengah genangan darah orang-orang tersayang dan reruntuhan fasilitas yang terbakar, mental serta jiwaku hancur sepenuhnya. Kekosongan yang begitu dingin menjalar dan meracuni seluruh dadaku.

Aku menengadah ke langit bungker yang gelap, melontarkan jeritan melengking yang membelah keheningan—jeritan dari hati yang telah mati.

``KENAPA?! KENAPA DUNIA INI SELALU MERAMPAS APA PUN YANG KUCINTAI?!''

Suaraku parau, tenggelam dalam isak tangis yang seakan meremukkan tulang rusuk. Air mata yang bercampur debu dan darah menetes ke wajah Kiren yang semakin pucat.

``Harapan$\dots$ harapan yang tak berarti!'' ratapku sambil dada terasa hancur berkeping-keping. ``Kalau memang ada harapan di dunia yang kejam ini$\dots$ lalu DI MANA DIA ketika orang-orang yang kucintai mati satu per satu?!''

Dua tahun telah berlalu. Usiaku kini masuk ke fase yang lebih dewasa, namun jiwaku mati dan membusuk di dalam tubuh yang kosong. Tanpa tujuan, tanpa tempat tinggal, tanpa kawan, dan tanpa sisa kebahagiaan sedikit pun. Aku bertahan hidup dari debu jalanan serta ceceran darah di arena pertarungan gelap. Setiap kali kepalaku berhasil menjatuhkan lawan, beberapa keping uang kotor kudapatkan untuk menyambung hidup—uang yang selalu habis tenggelam di meja-meja bar yang remang dan berbau busuk.

Aku sangat tahu bahwa air alkohol di masa krisis ini sungguh mahal, harganya melambung seperti setetes air bersih. Akan tetapi, candu melilit leherku begitu kencang. Rasa sakit dari kenangan pahit terlalu membakar jiwa, sampai aku lebih memilih meneguk racun pemabuk untuk mematikan rasa daripada membeli sesuap makanan.

12 November 2114.

Malam itu, langkah kakiku menyusuri lorong sektor dengan gontai, dikuasai kesadaran yang remang karena alkohol. Dari balik kegelapan lorong, pandanganku menangkap dua pria dewasa tengah memalak seorang anak kecil yang mungkin baru berusia sepuluh sampai tiga belas tahun. Awalnya, aku tidak peduli. Dunia ini memang sudah membusuk, dan empati di dadaku telah mati bersama Kiren dan Reymi.

Namun, ketika anak kecil itu berusaha melarikan diri, salah satu pria menyambar botol kaca tebal dan menghantamkannya keras ke kepala anak tersebut.

Prang!

Bunyi pecahannya memicu ledakan hebat di dalam kepalaku. Memori kelam saat Reymi dibantai dengan brutal oleh Pamanku tiba-tiba menyeruak membasahi pelipis. Trauma masa lalu berubah menjadi murka liar yang tak tertahankan.

Tanpa sadar, aku berlari menerjang. Pukulan mentah berkecepatan tinggi menghantam tepat pelipis pria pertama. Tengkoraknya menggelegar menghantam dinding beton sampai retak, meneteskan darah segar. Temannya yang panik berusaha membalas, melayangkan tinju yang mengenai punggung bawahku.

Rasa sakit itu justru menghanguskan sisa kemanusiaan yang kumiliki. Amarah yang berkobar-kobar menjadikanku sosok iblis. Aku menyerang pria kedua, mencengkeram lehernya, lalu membenturkan wajahnya berulang kali ke lantai semen tanpa ampun hingga jeritan kesakitannya lenyap ditelan malam. Anak kecil yang berusaha kuselamatkan justru melangkah mundur, menjerit ketakutan melihat kebrutalanku, lalu berlari kencang meninggalkanku sendirian bersama dua mayat yang tak lagi bernyawa.

Usai pembantaian itu, aku kembali meneguk sisa alkohol di genggamanku, melangkah tanpa tujuan menyusuri sektor demi sektor dalam kondisi mabuk berat.

Di sebuah sudut koridor berdebu, mata buramku menangkap sebuah poster lusuh yang tertempel di dinding metal. Sebuah turnamen pertarungan jalanan resmi yang diadakan oleh pemerintah—\textit{Street Champion}. Hadiah uang yang ditawarkan terlampau besar untuk diabaikan.

Selama ini aku memang kerap menjuarai laga pertarungan jalanan biasa, tetapi kali ini taruhannya berbeda. Pertarungan ini menyedot para petarung paling haus darah dari berbagai sektor. Dengan deru racun alkohol yang masih membakar darah, trauma yang terus merobek otak, dan himpitan beban kata harapan yang membusuk di dadaku, aku memantapkan langkah menuju arena neraka itu.

\jedahias
